\documentclass[aspectratio=169]{beamer}
\usetheme{Madrid}
\usecolortheme{default}
\usepackage[utf8]{inputenc}
\usepackage{amsmath}
\usepackage{amsfonts}
\usepackage{amssymb}

\setbeamertemplate{navigation symbols}{}
\setbeamerfont{block body}{size=\small}
\setbeamerfont{alertblock body}{size=\small}

\title{CRISC Domain B -- 100 Exam-Style Questions (Part 2)}
\subtitle{IT Risk Assessment -- Fresh Scenarios (Answer on the Next Slide)}
\author{CRISC Mastery Series}
\date{\today}

% ------------------------------------------------------------
% \hq{number}{topic}{question}{options}{letter}{explanation}{trap}
% Creates one question frame and one explanation frame.
% ------------------------------------------------------------
\newcommand{\hq}[7]{%
	\begin{frame}{Q#1: #2}
		\begin{block}{Question}
			#3
		\end{block}
		\begin{enumerate}[(A)]
			#4
		\end{enumerate}
	\end{frame}
	\begin{frame}{Q#1: Explanation}
		\begin{block}{Correct Answer: #5}
			\begin{itemize}
				#6
			\end{itemize}
		\end{block}
		\begin{alertblock}{Critical Trap}
			#7
		\end{alertblock}
	\end{frame}
}

\begin{document}
	
	\begin{frame}\titlepage\end{frame}
	
	\begin{frame}{Instructions}
		\begin{alertblock}{How to Use This Deck}
			Think like a risk practitioner: business objectives first, owners decide, assess before acting.\\
			Read for BEST, FIRST and MOST. Answer on the slide after each question.
		\end{alertblock}
	\end{frame}
	
	\begin{frame}{Outline}
		\tableofcontents
	\end{frame}
	
	% ============================================================
	\section{Risk Identification}
	% ============================================================
	
	\hq{1}{Strategic Alignment}
	{A risk practitioner is asked to help define the IT risk universe for the year. What should they do FIRST?}
	{\item List every known technical vulnerability
		\item Obtain the enterprise\textquotesingle s strategic objectives and business plans
		\item Benchmark against peers
		\item Issue a risk survey to all staff}
	{B}
	{\item The risk universe must support enterprise objectives
		\item Strategy defines what matters and what does not
		\item Peer benchmarks and surveys are supporting inputs, not the starting point
		\item Vulnerability lists are tactical, not strategic}
	{Start from enterprise strategy, not technology.}
	
	\hq{2}{Risk Appetite Input}
	{Which input is MOST important to the risk practitioner when scoping the annual IT risk assessment?}
	{\item The number of outstanding tickets
		\item The IT budget
		\item The enterprise risk appetite and prior risk profile
		\item The current technology stack}
	{C}
	{\item Appetite frames acceptable exposure
		\item The risk profile shows where you stand
		\item Tickets and stack are operational context, not scope drivers}
	{Scope should reflect appetite and profile.}
	
	\hq{3}{Business Process First}
	{When identifying IT risks for a new service, the BEST starting point is:}
	{\item The servers hosting the service
		\item The network diagrams
		\item The underlying codebase
		\item The business process the service supports and its objectives}
	{D}
	{\item Risk is about the effect on the business process
		\item Technology is the delivery mechanism
		\item Process objectives anchor relevance}
	{Understand the process before the plumbing.}
	
	\hq{4}{Third-Party Data}
	{The risk practitioner needs to identify risks from a new SaaS provider handling customer data. What should be requested FIRST?}
	{\item A copy of the vendor\textquotesingle s marketing brochure
		\item The vendor\textquotesingle s pricing list
		\item The vendor\textquotesingle s assurance report and data flow information
		\item The vendor\textquotesingle s org chart}
	{C}
	{\item Assurance reports show control posture
		\item Data flows show scope and location
		\item Marketing and pricing do not address risk}
	{Ask for assurance and data flow, not sales material.}
	
	\hq{5}{Change in Environment}
	{A new market expansion introduces new regulatory requirements. What should the risk practitioner do FIRST?}
	{\item Hire local staff
		\item Determine applicability and conduct a focused risk assessment
		\item Launch immediately and assess later
		\item Add insurance}
	{B}
	{\item Applicability frames the exposure
		\item A focused assessment identifies gaps
		\item Launching first increases exposure}
	{Assess before you expand.}
	
	\hq{6}{Scenario Source}
	{Which source BEST provides realistic risk scenarios for a mature organization?}
	{\item Anonymous suggestions only
		\item Generic industry checklists
		\item Lessons from past incidents plus intelligence on emerging threats
		\item The newest vendor whitepaper}
	{C}
	{\item Past incidents reveal real failure modes
		\item Emerging threat intelligence anticipates new risks
		\item Checklists and whitepapers are supporting material}
	{Combine experience with forward-looking intelligence.}
	
	\hq{7}{Loss Event Categories}
	{When building a loss event taxonomy, the practitioner should focus primarily on:}
	{\item Vendor names
		\item Loss categories that matter to the business such as confidentiality, integrity, availability and compliance
		\item Individual system owners
		\item IT budget lines}
	{B}
	{\item CIA plus compliance covers major impact types
		\item A consistent taxonomy supports aggregation
		\item Names and budgets are irrelevant to categorization}
	{Taxonomy must reflect business impact types.}
	
	\hq{8}{Interview Bias}
	{An interviewer asks leading questions about past incidents. What is the MAIN risk?}
	{\item The interview takes too long
		\item Responses may be biased toward the interviewer\textquotesingle s expectations
		\item The interviewee may refuse to answer
		\item The notes are hard to read}
	{B}
	{\item Leading questions steer answers
		\item Bias reduces reliability
		\item Use neutral, open questions and structure}
	{Neutral framing improves reliability.}
	
	\hq{9}{Risk Owner Identification}
	{During risk identification, what is MOST helpful in defining who will own a risk?}
	{\item The IT department head
		\item The external auditor
		\item The process owner who is accountable for the affected business outcome
		\item The risk practitioner}
	{C}
	{\item Ownership belongs to the person accountable for the outcome
		\item Process owners have the authority to act
		\item Practitioners and auditors advise}
	{Ownership follows the business outcome.}
	
	\hq{10}{Risk Workshop Output}
	{Which output is MOST valuable from a risk identification workshop?}
	{\item A long list of vulnerabilities
		\item A signed project charter
		\item Agreed risk statements with owners and business context
		\item A list of control costs}
	{C}
	{\item Risk statements link cause, event, and business impact
		\item Owners make them actionable
		\item Vulnerability lists are technical inputs}
	{Risk statements with owners are actionable.}
	
	\hq{11}{Asset Criticality}
	{When prioritizing asset-level risk identification, which asset should be reviewed first?}
	{\item The asset with the longest installed lifespan
		\item The asset with the highest business process criticality
		\item The newest asset
		\item The asset with the most users}
	{B}
	{\item Criticality reflects business consequence
		\item Older and newer assets come next based on criticality
		\item User count alone is not impact}
	{Criticality first, technology second.}
	
	\hq{12}{Data Classification}
	{What is the PRIMARY reason to involve data classification in risk identification?}
	{\item To reduce storage costs
		\item To determine which data needs protection and how much
		\item To make audits quicker
		\item To comply with HR policy}
	{B}
	{\item Classification drives control intensity
		\item Higher sensitivity means higher exposure
		\item Audits and costs are secondary effects}
	{Classification sets the protection bar.}
	
	\hq{13}{Shadow AI}
	{Employees use public generative AI tools with sensitive data. What should the practitioner do FIRST?}
	{\item Block all AI use
		\item Assess the data being exposed and the business impact
		\item Publicly discipline the users
		\item Accept the risk}
	{B}
	{\item Exposure type and volume determine severity
		\item Blocking may disrupt legitimate needs
		\item Acceptance requires assessment and authority}
	{Assess the exposure before reacting.}
	
	\hq{14}{Process Changes}
	{A core process is being redesigned. Why does the risk practitioner need to reassess risk?}
	{\item To add more meetings
		\item Because process changes can introduce new risks and retire old ones
		\item Because policies require a new form
		\item Because IT requested it}
	{B}
	{\item Changed processes change risk drivers
		\item New risks can emerge from design choices
		\item Old risks may be removed or replaced}
	{Redesign changes the risk landscape.}
	
	\hq{15}{External Threat Landscape}
	{Which external input MOST influences risk identification for internet-facing services?}
	{\item Peer benchmarking surveys
		\item Current threat intelligence and vulnerability disclosures relevant to the service
		\item Internal HR hiring plans
		\item Historic marketing reports}
	{B}
	{\item Threat intelligence reflects current attacker behavior
		\item Disclosures expose technical exposure
		\item Peer and HR inputs are not technical threat data}
	{Use current threat data for internet-facing risk.}
	
	% ============================================================
	\section{Risk Analysis}
	% ============================================================
	
	\hq{16}{Confidence Interval}
	{A quantitative analysis yields \$1.2M with a 90\% confidence interval of \$400K--\$2.0M. What does this mean?}
	{\item The risk is exactly \$1.2M
		\item There is a 90\% chance the true loss is between \$400K and \$2.0M
		\item The analysis is invalid
		\item The loss will not exceed \$2.0M}
	{B}
	{\item Confidence intervals reflect uncertainty
		\item The point estimate is not exact
		\item Wide ranges signal high uncertainty}
	{Report the range, not just the point.}
	
	\hq{17}{Probability Distribution}
	{Which distribution is MOST appropriate for modeling the number of independent incidents in a period?}
	{\item Uniform
		\item Poisson
		\item Normal
		\item Deterministic}
	{B}
	{\item Poisson models counts of independent events
		\item Uniform assumes equal likelihood
		\item Normal suits continuous measures
		\item Deterministic ignores uncertainty}
	{Choose the model that matches the data.}
	
	\hq{18}{Monte Carlo}
	{Monte Carlo simulation is MOST useful when:}
	{\item Only one outcome is possible
		\item There are many interacting uncertain variables
		\item The data is fully known
		\item No computation tools are available}
	{B}
	{\item Simulation captures interactions of uncertainties
		\item Single-outcome and fully known cases do not need it
		\item Computation tools are required}
	{Simulate when uncertainty is complex.}
	
	\hq{19}{Data Normalization}
	{Two business units report losses in different units and currencies. What should the analyst do FIRST?}
	{\item Average the raw numbers
		\item Ignore the smaller unit
		\item Normalize units and currency to enable comparison
		\item Use only the larger unit}
	{C}
	{\item Mixed units produce invalid comparisons
		\item Normalization enables aggregation
		\item Ignoring data biases results}
	{Normalize before you compare.}
	
	\hq{20}{Subject Matter Experts}
	{Which situation reduces the value of subject matter expert judgment?}
	{\item Experts are drawn from multiple functions
		\item Experts provide documented rationale
		\item Experts are interviewed independently
		\item Experts share a single worldview and are not diverse}
	{D}
	{\item Diverse perspectives reduce blind spots
		\item Independent inputs reduce groupthink
		\item Documented rationale improves auditability
		\item Homogeneous experts amplify bias}
	{Diversity beats sameness in judgment.}
	
	\hq{21}{Benchmarking Limits}
	{Benchmarking with peer breach data should be used primarily to:}
	{\item Replace internal assessment
		\item Inform context, not to substitute for internal analysis
		\item Set the final risk score
		\item Avoid collecting internal data}
	{B}
	{\item Peer data may not fit your environment
		\item Internal data reflects your controls and threats
		\item Benchmarks inform, not replace}
	{Benchmarks inform; they do not decide.}
	
	\hq{22}{Risk Score Aggregation}
	{The enterprise risk profile contains many small IT risks. What is the MOST likely mistake in aggregation?}
	{\item Underestimating correlated exposure
		\item Using a single scoring scale
		\item Documenting all risks
		\item Involving business owners}
	{A}
	{\item Related risks compound beyond individual ratings
		\item Correlated causes increase total exposure
		\item Documentation and participation are good practices}
	{Watch correlation when aggregating.}
	
	\hq{23}{Risk Sensitivity}
	{A small change in probability assumption changes the rating from Medium to High. What does this indicate?}
	{\item The rating is stable
		\item The analysis is highly sensitive to that assumption and should be validated
		\item The risk should be ignored
		\item The probability is wrong}
	{B}
	{\item Sensitivity signals fragility of results
		\item Validate the assumption before acting
		\item Ignoring the signal leads to misallocation}
	{Test sensitive assumptions before decisions.}
	
	\hq{24}{ALE Trap}
	{An asset is worth \$1M, the exposure factor is 20\%, and ARO is 3. What is ALE?}
	{\item \$60K
		\item \$2.4M
		\item \$200K
		\item \$600K}
	{D}
	{\item SLE $= 1M \times 0.2 = \$200K$
		\item ALE $= 200K \times 3 = \$600K$
		\item \$200K is only SLE}
	{ALE is SLE multiplied by ARO.}
	
	\hq{25}{Expected Versus Extreme}
	{The expected value of a risk is low but the worst-case outcome is severe. What should the risk manager present?}
	{\item Only the expected value
		\item Only the worst case
		\item Both expected value and worst case, with a recommendation
		\item Neither, to avoid confusion}
	{C}
	{\item Decisions need a full picture
		\item Expected value guides normal planning
		\item Worst case informs contingency and appetite}
	{Present both expected and worst case.}
	
	\hq{26}{Model Assumptions}
	{Before approving a risk analysis, leadership asks about assumptions. Why is this important?}
	{\item To increase the length of the report
		\item To test whether the conclusions hold if assumptions change
		\item To reduce the budget
		\item To avoid the risk}
	{B}
	{\item Assumptions shape results
		\item Leadership needs to judge reliability
		\item Documenting assumptions supports governance}
	{Assumptions drive conclusions.}
	
	\hq{27}{Cluster of Risks}
	{You discover ten risks tied to a single outdated platform. What is the BEST analytical view?}
	{\item Treat them as independent
		\item Ignore the platform
		\item Look at them as a cluster and assess platform-level exposure
		\item Remove the lowest-rated ones}
	{C}
	{\item Common causes create correlated risk
		\item A platform-level view reveals total exposure
		\item Independence assumptions understate total risk}
	{Common causes mean combined exposure.}
	
	\hq{28}{Data Age Impact}
	{Threat intelligence is fresh, but internal control data is two years old. Which combination is BEST for the analysis?}
	{\item Fresh external threat data combined with refreshed internal control data
		\item Fresh external data only
		\item Old internal data only
		\item Neither}
	{A}
	{\item Fresh external data is current risk
		\item Stale internal data misrepresents controls
		\item Both must be current for reliable results}
	{Refresh internal data before using it.}
	
	\hq{29}{Risk Register Aggregate}
	{The number of high-rated IT risks in the register has doubled in one quarter. What is the FIRST interpretation to verify?}
	{\item The business is definitely more exposed
		\item A change in scoring methodology or assessment quality
		\item The controls stopped working
		\item The register is wrong}
	{B}
	{\item Methodology changes drive rating shifts
		\item Verify before concluding exposure increased
		\item Controls failure would be a specific finding}
	{Check methodology before concluding trend.}
	
	\hq{30}{Risk Rating Consistency}
	{Two locations use the same methodology but differ in rating severity. What is the BEST remedy?}
	{\item Standardized calibration workshops with real examples
		\item Report both ratings
		\item Average the ratings
		\item Use only the higher rating}
	{A}
	{\item Calibration aligns interpretation
		\item Examples anchor scoring
		\item Averages and higher-only hide inconsistency}
	{Calibrate with examples.}
	
	% ============================================================
	\section{Risk Evaluation and Appetite}
	% ============================================================
	
	\hq{31}{Appetite Translation}
	{The board expresses appetite as protecting brand reputation. What should the practitioner do NEXT?}
	{\item Ignore the statement
		\item Translate the appetite into measurable thresholds for IT risk
		\item Buy insurance
		\item Freeze the risk register}
	{B}
	{\item Appetite must be operationalized
		\item Measurable thresholds enable comparison
		\item Protection actions follow thresholds}
	{Translate appetite into measurable thresholds.}
	
	\hq{32}{Risk Owner Escalation}
	{The risk owner wants to accept a risk well above appetite. What should the practitioner do?}
	{\item Accept it, the owner is accountable
		\item Escalate according to governance for higher authority
		\item Reduce the rating
		\item Delete the risk entry}
	{B}
	{\item Accepting above appetite requires higher authority
		\item The practitioner must follow escalation policy
		\item Ratings must not be manipulated}
	{Acceptance above appetite is not the owner’s call.}
	
	\hq{33}{Risk Tolerance Example}
	{Which is the BEST example of risk tolerance?}
	{\item Loss of up to 3 hours of transaction data before escalating
		\item Zero incidents allowed
		\item No data loss ever
		\item Unbounded loss}
	{A}
	{\item Tolerance is a defined acceptable variation
		\item Zero and unbounded are not realistic thresholds
		\item Specific measure supports monitoring}
	{Tolerance is a defined, measurable variation.}
	
	\hq{34}{Post-Treatment Above Appetite}
	{Residual risk after the planned controls is still above appetite. What is the BEST course?}
	{\item Close the risk
		\item Escalate for additional treatment or formal acceptance
		\item Lower the rating
		\item Ignore it}
	{B}
	{\item Treatment or acceptance must be decided
		\item The owner or higher authority decides
		\item Quiet closure hides exposure}
	{Above appetite needs an explicit decision.}
	
	\hq{35}{Risk Acceptance Criteria}
	{What should accompany a formal risk acceptance?}
	{\item A signature only
		\item Documentation of impact, rationale, duration and approval authority
		\item A verbal note
		\item No documentation}
	{B}
	{\item Documented rationale supports accountability
		\item Duration and authority must be clear
		\item Undocumented acceptance is indefensible}
	{Acceptance needs evidence.}
	
	\hq{36}{Appetite Monitoring}
	{What is the PRIMARY purpose of linking KRIs to risk appetite?}
	{\item To reduce reporting
		\item To detect when exposure is approaching the acceptable threshold
		\item To justify budget increases
		\item To satisfy auditors}
	{B}
	{\item Appetite is a threshold
		\item KRIs detect drift
		\item Reporting is a means, not the goal}
	{KRIs are appetite monitors.}
	
	\hq{37}{Capacity Breach}
	{Exposure is approaching risk capacity. What should be prioritized?}
	{\item Cost reduction
		\item Immediate executive action to reduce exposure
		\item Policy updates
		\item Staff training}
	{B}
	{\item Capacity is existential
		\item Immediate action is required
		\item Other measures follow}
	{Capacity signals emergency action.}
	
	\hq{38}{Appetite Versus Mission}
	{A new business line exceeds appetite but aligns with mission. What is the BEST guidance?}
	{\item Proceed without discussion
		\item Escalate to leadership to align appetite and opportunity
		\item Abandon the line
		\item Ignore appetite}
	{B}
	{\item Appetite may need adjustment
		\item Leadership balances mission and risk
		\item Silent proceeding is not governance}
	{Opportunity requires re-aligning appetite.}
	
	\hq{39}{Inconsistent Impact}
	{The business wants to reduce impact ratings to fit the budget. What is the practitioner’s BEST response?}
	{\item Reduce ratings to align with budget
		\item Present evidence and independently validated impact assessments
		\item Leave it to IT
		\item Approve it}
	{B}
	{\item Impact must be based on evidence
		\item Independence prevents manipulation
		\item Approving manipulated ratings undermines governance}
	{Ratings reflect evidence, not budgets.}
	
	\hq{40}{Accepted Risk Monitoring}
	{An accepted risk is now trending toward its tolerance limit. What should happen?}
	{\item Remove the risk
		\item Re-open the decision with the owner and consider new treatment
		\item Ignore the trend
		\item Cancel the KRI}
	{B}
	{\item Acceptance is not permanent
		\item Trends trigger re-evaluation
		\item Ignoring trends misses warnings}
	{Acceptance must be re-affirmed, not assumed.}
	
	% ============================================================
	\section{Risk Response}
	% ============================================================
	
	\hq{41}{Transfer Limitations}
	{Which risk can insurance NOT transfer?}
	{\item Reputational loss from a public breach
		\item Financial cost of recovery
		\item Legal defense costs
		\item Regulatory fines where insurable}
	{A}
	{\item Reputation is not insurable
		\item Financial and legal costs are insurable
		\item Some fines are insurable depending on jurisdiction}
	{Reputation cannot be transferred.}
	
	\hq{42}{Justified Investment}
	{An annual ALE of \$500K is addressed by a control costing \$200K. Which statement is MOST accurate?}
	{\item The control is not justified
		\item The control may be justified pending effectiveness validation
		\item The control is always unnecessary
		\item No analysis is required}
	{B}
	{\item Cost is less than the ALE
		\item Effectiveness must be validated
		\item Justification is supported but conditional}
	{Cost-benefit supports it; validate effectiveness.}
	
	\hq{43}{Outsourcing Oversight}
	{After outsourcing a process, what oversight should continue?}
	{\item None, the vendor is responsible
		\item Monitoring of the vendor’s performance and assurance evidence
		\item Only annual meetings
		\item Only contract renewal}
	{B}
	{\item Outsourcing transfers tasks, not accountability
		\item Ongoing monitoring catches drift
		\item Annual meetings are not enough}
	{Monitor the vendor continuously.}
	
	\hq{44}{Response Alignment}
	{Which response is BEST aligned with a risk appetite that tolerates minor outages but not data loss?}
	{\item Reduce availability controls and focus on backups and integrity controls
		\item Remove backups
		\item Focus only on uptime marketing
		\item Do nothing}
	{A}
	{\item The appetite differentiates availability and data loss
		\item Integrity and backup controls address data loss
		\item Availability controls can be lighter}
	{Fit the control portfolio to appetite specifics.}
	
	\hq{45}{Avoidance Decision}
	{When is avoidance the MOST appropriate response?}
	{\item When the risk can be reduced to acceptable levels cheaply
		\item When the risk cannot be treated to acceptable levels and the activity is not essential
		\item When the risk is low
		\item When insurance is available}
	{B}
	{\item Avoidance stops the activity
		\item It suits unacceptable, inessential risks
		\item Other options are cheaper when viable}
	{Avoid only when the activity is truly not needed.}
	
	\hq{46}{Preventive Example}
	{Which control is PURELY preventive?}
	{\item Log monitoring
		\item Backup restoration
		\item Encryption at rest
		\item Incident reporting}
	{C}
	{\item Encryption prevents reading even if data is exposed
		\item Log monitoring is detective
		\item Backup restoration is corrective}
	{Prevention stops exposure before it matters.}
	
	\hq{47}{Detective Timing}
	{A detective control fails silently for six months. What is the MOST likely consequence?}
	{\item No consequence
		\item Incidents may go undetected and their impact will grow
		\item The risk is eliminated
		\item The threat disappears}
	{B}
	{\item Detection failure delays response
		\item Impact grows
		\item Visibility of failures is critical}
	{Detective controls fail silently.}
	
	\hq{48}{Compensating Approval}
	{Who should approve a compensating control?}
	{\item The vendor
		\item The control owner only
		\item The risk owner, based on evidence that it addresses the same risk
		\item The IT helpdesk}
	{C}
	{\item Compensating controls must reduce the same risk
		\item The risk owner decides whether it is sufficient
		\item The control owner documents operation}
	{The risk owner approves compensation.}
	
	\hq{49}{Residual Ownership}
	{Who is accountable for residual risk after implementation?}
	{\item The vendor
		\item The project manager
		\item The IT department
		\item The risk owner}
	{D}
	{\item Residual risk remains with the business owner
		\item The project ends, but the risk continues
		\item Vendors do not hold business accountability}
	{Residual risk stays with the owner.}
	
	\hq{50}{Response Effectiveness}
	{Which test MOST reliably confirms a risk response is effective?}
	{\item Reviewing the design document
		\item A verified post-implementation review with evidence
		\item Verbal assurance from the team
		\item A project closure report}
	{B}
	{\item Evidence confirms real effect
		\item Design reviews do not prove operation
		\item Verbal and closure reports may lack rigor}
	{Verify with evidence.}
	
	\hq{51}{Documented Accept}
	{A risk owner refuses to sign the acceptance. What should the practitioner do?}
	{\item Proceed as if accepted
		\item Document the refusal and escalate
		\item Forge the acceptance
		\item Delete the risk}
	{B}
	{\item Refusal indicates unresolved concerns
		\item Escalation is required
		\item Undocumented or forged acceptance is improper}
	{Refusal is a governance event.}
	
	\hq{52}{Vendor Assurance}
	{The vendor will not share third-party assurance reports but permits an audit. What is the BEST path?}
	{\item Conduct the audit or accept additional compensating controls
		\item Ignore the vendor
		\item Sign anyway
		\item Skip the contract}
	{A}
	{\item Direct audit provides assurance
		\item Compensating controls may be acceptable
		\item Skipping assurance leaves exposure unknown}
	{Audit or compensate, do not skip.}
	
	\hq{53}{Exception Cadence}
	{A policy exception is approved for one year. What should happen at year-end?}
	{\item Renew automatically
		\item Validate the need and controls, then decide
		\item Delete the exception
		\item Escalate to the board}
	{B}
	{\item Circumstances may have changed
		\item Validation prevents stale exceptions
		\item Automatic renewal is not governance}
	{Exceptions must be re-validated.}
	
	\hq{54}{Ineffective Response}
	{A control is not delivering the expected risk reduction. What should happen FIRST?}
	{\item Remove the control
		\item Diagnose whether it is a design or operating issue
		\item Increase the risk rating
		\item Cancel audits}
	{B}
	{\item The diagnosis determines the fix
		\item Design and operating failures need different remedies
		\item Removing controls may increase exposure}
	{Diagnose before changing.}
	
	\hq{55}{Justified Response}
	{Which condition most strongly justifies mitigation?}
	{\item Risk exceeds appetite and a cost-effective control exists
		\item Risk is low
		\item Risk has never occurred
		\item Insurance is available}
	{A}
	{\item Above-appetite risk needs action
		\item Cost-effectiveness supports the investment
		\item Low and non-occurring risks do not force mitigation}
	{Mitigate when above appetite and cost-effective.}
	
	% ============================================================
	\section{Controls}
	% ============================================================
	
	\hq{56}{Deficiency Type}
	{A control is well designed but the tool is misconfigured. This is a:}
	{\item Design deficiency
		\item Operating deficiency
		\item Governance failure
		\item Scope failure}
	{B}
	{\item Design is sound
		\item Operation is defective
		\item Remediation is configuration}
	{Configuration failures are operating issues.}
	
	\hq{57}{Layered Controls}
	{Which combination BEST represents defense in depth?}
	{\item Multiple controls sharing the same weakness
		\item One control only
		\item Diverse preventive, detective and corrective controls with independent failure modes
		\item Duplicate tools from the same vendor}
	{C}
	{\item Independence reduces common-mode failures
		\item Diversity matters more than quantity
		\item Duplicates may share the same flaw}
	{Layers need independence.}
	
	\hq{58}{Control Ownership}
	{Which statement BEST describes control ownership?}
	{\item The owner is responsible for design, operation and evidence
		\item The owner only pays for the control
		\item The owner audits the control
		\item The owner approves risk appetite}
	{A}
	{\item Ownership includes operation and evidence
		\item Auditing is independent
		\item Appetite belongs to leadership}
	{Control ownership includes operation and evidence.}
	
	\hq{59}{Restoration Time}
	{A backup restores data but takes 24 hours against a 4-hour RTO. This is:}
	{\item A corrective control that meets its objective
		\item A corrective control that does not meet the recovery objective
		\item A preventive control
		\item A detective control}
	{B}
	{\item Restoration is corrective
		\item The RTO is missed
		\item The control is failing a key requirement}
	{Corrective controls must meet recovery targets.}
	
	\hq{60}{Independence}
	{Which tester provides the MOST objective assurance?}
	{\item The control owner
		\item The project manager
		\item The vendor
		\item An independent assurance function not responsible for the control}
	{D}
	{\item Independence prevents bias
		\item Responsibility reduces objectivity
		\item Vendors have commercial interests}
	{Independence ensures objectivity.}
	
	\hq{61}{Gap Handling}
	{A control gap is found. What should be done FIRST?}
	{\item Report to the board
		\item Assess the exposure caused by the gap
		\item Implement a tool
		\item Change policy}
	{B}
	{\item Exposure determines urgency
		\item Urgent risks get immediate action
		\item Lower exposure allows prioritization}
	{Assess exposure before fixing.}
	
	\hq{62}{Automation Risk}
	{Which risk does automation introduce?}
	{\item No risk
		\item Consistent operation
		\item Silent failure if monitoring is absent
		\item Lower cost}
	{C}
	{\item Automated controls can fail silently
		\item Monitoring is required
		\item Consistency is a benefit, not a risk}
	{Automate, but monitor silently.}
	
	\hq{63}{Fit to Appetite}
	{How should control selection relate to appetite?}
	{\item Controls should reduce risk to within appetite at justified cost
		\item Any control is fine
		\item Controls should exceed appetite
		\item Cost is the only factor}
	{A}
	{\item Appetite drives the target level
		\item Cost must be justified
		\item Appetite alignment makes controls proportionate}
	{Match controls to appetite and cost.}
	
	\hq{64}{Redundant Controls}
	{Two independent controls cover the same risk. What should the practitioner do?}
	{\item Remove one immediately
		\item Evaluate whether redundancy supports defense in depth
		\item Ignore both
		\item Increase the risk rating}
	{B}
	{\item Redundancy can be desirable
		\item Evaluate effectiveness and cost
		\item Immediate removal may increase exposure}
	{Justify redundancy, do not assume waste.}
	
	\hq{65}{CSA Value}
	{Where does control self-assessment provide the GREATEST value?}
	{\item When owners understand and improve their own controls
		\item When there is no external validation
		\item When used to replace audits
		\item When owners rate themselves highly}
	{A}
	{\item CSA drives ownership and improvement
		\item External validation adds assurance
		\item CSA complements audits}
	{CSA drives ownership and improvement.}
	
	% ============================================================
	\section{Monitoring and Reporting}
	% ============================================================
	
	\hq{66}{Leading Indicator}
	{Which indicator is most useful as a leading KRI for a ransomware risk?}
	{\item Historical ransomware incidents
		\item Patch latency and backup restore success trends
		\item Number of unpatched servers last year
		\item Length of the last downtime}
	{B}
	{\item Patch latency and restore success predict future impact
		\item Historical incidents are lagging
		\item Prior-year counts are lagging
		\item Downtime length is a lagging outcome}
	{Leading KRIs predict; lagging ones record.}
	
	\hq{67}{Signal vs Noise}
	{A KRI routinely alerts during business-critical peak hours. What is the BEST fix?}
	{\item Ignore it
		\item Tune thresholds and correlate with business context
		\item Disable it
		\item Increase the risk rating}
	{B}
	{\item Context and thresholds determine false positives
		\item Tuning improves signal
		\item Disabling removes warning}
	{Tune, do not disable.}
	
	\hq{68}{Audience Fit}
	{A regulator requests detailed evidence of monitoring. Which report BEST fits?}
	{\item A board-level summary
		\item A concise executive dashboard
		\item A detailed, evidence-based monitoring report
		\item A marketing brochure}
	{C}
	{\item Regulators need evidence and detail
		\item Boards need summaries
		\item Evidence supports compliance}
	{Fit the detail to the audience.}
	
	\hq{69}{Register Trigger}
	{Which event should NOT automatically trigger a risk register update?}
	{\item A significant incident
		\item A new threat intelligence update relevant to an asset
		\item A change in the risk owner
		\item A routine staff birthday}
	{D}
	{\item Incidents change exposures
		\item Threat updates change likelihood
		\item Ownership changes affect accountability}
	{Not everything updates the register.}
	
	\hq{70}{Real-Time Alert}
	{Which scenario MOST benefits from real-time monitoring?}
	{\item Annual compliance reporting
		\item Controls protecting high-value transactions
		\item Long-term strategy reviews
		\item Training attendance}
	{B}
	{\item High-value transactions require timely detection
		\item Annual and strategic activities are periodic
		\item Training is monitored periodically}
	{Real-time where impact is immediate.}
	
	\hq{71}{KPI Versus KRI}
	{Which metric is a KPI rather than a KRI?}
	{\item Percentage of phishing emails that bypassed filters
		\item Number of vulnerabilities with active exploits
		\item Time to respond to incidents against service targets
		\item Amount of data exposed}
	{C}
	{\item Response time is performance against a target
		\item Phishing and exploit-based metrics measure risk
		\item Exposure volume is a risk outcome}
	{KPIs measure performance against targets.}
	
	\hq{72}{Overload}
	{The risk team receives 1,000 alerts per day and responds to none. What is the BEST fix?}
	{\item Add more alerts
		\item Reduce alerts to a prioritized, actionable set with response procedures
		\item Ignore the alerts
		\item Turn off the tool}
	{B}
	{\item Prioritization improves response
		\item Alerts need owners and procedures
		\item Disabling removes visibility}
	{Make alerts actionable and prioritized.}
	
	\hq{73}{Stable Indicators}
	{All indicators have been stable and green for three years. What is the BEST action?}
	{\item Change nothing
		\item Challenge whether the metrics remain relevant and sensitive
		\item Reduce the tool budget
		\item Increase thresholds}
	{B}
	{\item Relevance must be validated over time
		\item Insensitivity hides risk
		\item Tuning thresholds upward reduces visibility}
	{Challenge stable indicators.}
	
	\hq{74}{Board Message}
	{Which format BEST conveys risk to the board?}
	{\item A raw KRI list
		\item A concise view of key exposures against appetite with trends
		\item A technical log
		\item A list of all vulnerabilities}
	{B}
	{\item Boards need decision-level insight
		\item Appetite and trends guide oversight
		\item Raw and technical details overwhelm}
	{Boards need clarity, not detail.}
	
	\hq{75}{Missed Signal}
	{A breach occurred but no KRI flagged it. What is the BEST response?}
	{\item Ignore it
		\item Review the monitoring design and consider new or adjusted KRIs
		\item Blame staff
		\item Retire the framework}
	{B}
	{\item The framework missed a signal
		\item Framework improvement is the goal
		\item Blaming individuals does not improve controls}
	{Missed signals mean improve the framework.}
	
	\hq{76}{Register Reviews}
	{Which cadence BEST supports register accuracy?}
	{\item Only on incidents
		\item Periodic reviews by owners with risk practitioner facilitation
		\item Never
		\item Only at year-end}
	{B}
	{\item Owners keep entries current
		\item Periodic reviews catch drift
		\item Incidents and year-end are triggers but not sufficient alone}
	{Periodic owner reviews keep registers accurate.}
	
	\hq{77}{Escalation Design}
	{What MOST improves the timeliness of risk escalation?}
	{\item Formal written reports only
		\item Predefined triggers and named escalation recipients
		\item Verbose minutes
		\item Long committee cycles}
	{B}
	{\item Triggers speed action
		\item Named recipients clarify accountability
		\item Written-only and long cycles slow escalation}
	{Predefine triggers and owners.}
	
	\hq{78}{KRI Reliability}
	{Which source MOST reduces KRI manipulation risk?}
	{\item Self-reported metrics only
		\item Automated or independent data feeds
		\item Verbal summaries
		\item Emailed spreadsheets}
	{B}
	{\item Automation reduces manipulation
		\item Independence strengthens trust
		\item Verbal and spreadsheet data is harder to verify}
	{Trust automated or independent data.}
	
	\hq{79}{Audit Integration}
	{How should audit findings integrate with risk monitoring?}
	{\item They replace KRIs
		\item They update risks and trigger reassessment where relevant
		\item They should be ignored
		\item They are only for compliance reporting}
	{B}
	{\item Findings reveal control weaknesses
		\item Findings may change likelihood and impact
		\item They add to, not replace, monitoring}
	{Audit findings feed the risk process.}
	
	\hq{80}{Trend Value}
	{What does an increasing KRI trend below threshold indicate?}
	{\item No concern
		\item Exposure may be rising and may breach threshold soon
		\item The KRI is wrong
		\item The risk is eliminated}
	{B}
	{\item Trends can precede breaches
		\item Early action is possible
		\item Ignoring trends reduces response time}
	{Rising trends deserve early action.}
	
	% ============================================================
	\section{BIA, Recovery and Third Parties}
	% ============================================================
	
	\hq{81}{BIA Owner}
	{Which group should OWN the business impact analysis?}
	{\item IT operations
		\item Internal audit
		\item Business process owners
		\item Vendors}
	{C}
	{\item Business owners understand impact
		\item IT provides recovery capability
		\item Auditors provide assurance}
	{The business owns the BIA.}
	
	\hq{82}{RTO Example}
	{Which statement BEST describes an RTO of 6 hours?}
	{\item Data can be lost for up to 6 hours
		\item The process should be restored within 6 hours of disruption
		\item Downtime is allowed up to 6 hours before consideration of action
		\item 6 hours is the maximum outage tolerated}
	{B}
	{\item RTO is a recovery target
		\item Data loss is RPO
		\item Tolerance limits are separate}
	{RTO is time to restore.}
	
	\hq{83}{RPO Example}
	{Which statement BEST describes an RPO of 1 hour?}
	{\item Recovery must occur within 1 hour
		\item No more than 1 hour of data can be lost
		\item Backups must run every 24 hours
		\item Downtime can be 1 hour}
	{B}
	{\item RPO is a data loss limit
		\item Backup frequency should match
		\item Downtime is RTO}
	{RPO is data loss tolerance.}
	
	\hq{84}{RTO and MTD}
	{Which relationship should hold for a critical system?}
	{\item RTO greater than MTD
		\item RTO equal to MTD
		\item RTO less than or equal to MTD
		\item No relationship}
	{C}
	{\item Recovery must occur within the tolerable limit
		\item RTO at or below MTD protects the business
		\item RTO beyond MTD leaves impact unaddressed}
	{RTO stays within MTD.}
	
	\hq{85}{BIA Prioritization}
	{Which criterion BEST drives BIA prioritization?}
	{\item Number of users
		\item Cost of the system
		\item Business and financial impact over time
		\item Vendor preference}
	{C}
	{\item Impact drives priorities
		\item Users and cost are inputs but not primary
		\item Vendors are not the business}
	{Impact over time determines priority.}
	
	\hq{86}{Backup Gap}
	{RPO is 4 hours; nightly backups run at midnight. What is the exposure?}
	{\item Up to 4 hours of data loss
		\item Up to 24 hours of data loss
		\item No data loss
		\item Up to 1 hour of data loss}
	{B}
	{\item Nightly backups allow up to 24 hours of loss
		\item The RPO is exceeded
		\item Backup design must change}
	{Match backup cadence to RPO.}
	
	\hq{87}{DR Validation}
	{A DR test reveals a 10-hour restore against a 4-hour RTO. What should be done NEXT?}
	{\item Accept the result
		\item Identify causes, update the plan and retest
		\item Raise the RTO automatically
		\item Cancel the test}
	{B}
	{\item RTO misses must be addressed
		\item Root causes and retests improve confidence
		\item Unilateral RTO change needs business approval}
	{Fix and retest, do not ignore.}
	
	\hq{88}{Vendor Risk Timing}
	{When should a third party be evaluated for the FIRST time?}
	{\item After an incident
		\item During contract negotiation and before onboarding
		\item Only at renewal
		\item Never}
	{B}
	{\item Pre-contract evaluation shapes terms
		\item Negotiation enables controls
		\item Onboarding without evaluation increases exposure}
	{Evaluate before onboarding.}
	
	\hq{89}{Contract Controls}
	{Which contract clause MOST supports ongoing third-party risk oversight?}
	{\item A pricing clause
		\item Right-to-audit and assurance reporting duties
		\item Logo usage
		\item Coffee supply at meetings}
	{B}
	{\item Oversight needs access to evidence
		\item Reporting duties provide ongoing visibility
		\item Commercial items do not address risk}
	{Oversight depends on contract rights.}
	
	\hq{90}{Concentration}
	{An organization relies on one cloud provider for all critical services. What is the BEST action?}
	{\item Increase reliance
		\item Assess concentration risk and develop contingency or exit strategy
		\item Ignore it
		\item Cancel the contract}
	{B}
	{\item Concentration is a risk
		\item Contingency reduces impact
		\item Ignoring the risk understates exposure}
	{Concentration needs an exit plan.}
	
	% ============================================================
	\section{Governance and Risk Culture}
	% ============================================================
	
	\hq{91}{Register Content}
	{Which entry set is MINIMUM for a credible risk register entry?}
	{\item Owner, risk statement, rating, response, review date
		\item Contact phone number and budget code
		\item Marketing description
		\item Vendor name only}
	{A}
	{\item Owners, statements and ratings enable management
		\item Review dates keep entries fresh
		\item Extraneous fields add noise}
	{Core fields enable governance.}
	
	\hq{92}{Owner Authority}
	{A risk owner lacks the authority to approve needed spending. What should happen?}
	{\item Approve without authority
		\item Realign ownership with someone who has the authority
		\item Ignore the risk
		\item Lower the rating}
	{B}
	{\item Ownership requires authority
		\item Misalignment causes ineffective action
		\item Ratings must not be manipulated}
	{Ownership requires the ability to act.}
	
	\hq{93}{Tone at the Top}
	{Which action MOST signals a strong risk culture?}
	{\item Leaders acknowledge risks and act on them
		\item Frequent communications only
		\item Detailed dashboards only
		\item Annual celebrations}
	{A}
	{\item Leadership behavior drives culture
		\item Communication and dashboards support, but do not replace, behavior
		\item Celebrations do not change behavior}
	{Behaviors shape culture.}
	
	\hq{94}{Line Roles}
	{In the three lines model, who provides independent assurance?}
	{\item The first line
		\item The second line
		\item The third line (internal audit)
		\item External vendors}
	{C}
	{\item First line owns risk
		\item Second line oversees and advises
		\item Third line provides independent assurance}
	{Independence sits in the third line.}
	
	\hq{95}{Practitioner Behavior}
	{Which action BEST reflects the risk practitioner’s independence?}
	{\item Approving the risk acceptance
		\item Advising the owner while avoiding the decision
		\item Choosing the vendor
		\item Operating the control}
	{B}
	{\item Independence means advising, not deciding
		\item Choosing vendors or operating controls reduces independence
		\item Owners decide}
	{Advise without deciding.}
	
	\hq{96}{Emerging Tech}
	{Before adopting a new technology, which governance step is MOST important?}
	{\item A pilot without assessment
		\item A proportionate risk assessment and ownership decision
		\item Immediate production rollout
		\item Waiting for incidents}
	{B}
	{\item Assessment informs decisions
		\item Proportionality keeps adoption practical
		\item Unassessed rollout increases exposure}
	{Assess before adoption.}
	
	\hq{97}{Frequency}
	{Which combination BEST defines risk assessment frequency?}
	{\item Annual only
		\item On incidents only
		\item Periodically and on significant change
		\item Only at project end}
	{C}
	{\item Periodic reviews catch drift
		\item Significant changes create new risks
		\item Annual is a minimum, not sufficient alone}
	{Periodic plus change-driven assessments.}
	
	\hq{98}{Executive Language}
	{Which phrase BEST supports executive decision-making?}
	{\item Three CVEs have a CVSS score above 9
		\item The payment process has a 30\% chance of extended outage within twelve months
		\item Patches were deployed to VLAN 12
		\item The SIEM correlates logs by rule set}
	{B}
	{\item Executives need business impact and likelihood
		\item Vulnerability and configuration detail belong to specialists
		\item Business-oriented language drives action}
	{Speak in business consequences.}
	
	\hq{99}{Stale Entries}
	{An owner has not reviewed their register entries for two years. What is the BEST response?}
	{\item Escalate stale entries and reset a review cadence
		\item Remove the owner
		\item Ignore them
		\item Increase ratings}
	{A}
	{\item Stale entries may be wrong
		\item Escalation encourages ownership
		\item A cadence prevents drift}
	{Stale registers misinform.}
	
	\hq{100}{Owner Assigned}
	{A new risk has no owner. What should the practitioner do?}
	{\item Assign it to IT
		\item Escalate to leadership to assign an accountable owner
		\item Leave it unowned
		\item Close it}
	{B}
	{\item Ownership must reflect accountability
		\item Leadership assigns owners
		\item Defaulting to IT misaligns accountability}
	{Unowned risks are unmanaged.}
	
	\begin{frame}{Exam Mindset Summary}
		\begin{block}{Remember}
			\begin{itemize}
				\item \textbf{FIRST:} understand context and assess before acting
				\item \textbf{Owners decide:} business owners accept risk; the practitioner advises
				\item \textbf{Appetite:} compare residual risk to appetite, then escalate or treat
				\item \textbf{Accountability} cannot be outsourced or transferred
				\item \textbf{Monitor:} validate data, trend it, and keep the register current
			\end{itemize}
		\end{block}
	\end{frame}
	
	\begin{frame}{Thank You}
		\begin{center}
			{\Huge \textbf{Good Luck on Your CRISC Exam!}}
		\end{center}
	\end{frame}
	
\end{document}